\documentclass[
	10pt,								% globale Schriftgröße
	parskip=half-,						% setzt Absatzabstand hoch
	paper=a4,							% Format
	english,ngerman,					% lädt Sprachpakete
	]{scrartcl}							% Dokumentenklasse

% //////////////////// Pakete laden ////////////////////
\usepackage[fleqn]{amsmath}
\usepackage[fleqn]{mathtools}
\usepackage{amssymb}			% mathematische symbole, für \ceckmarks
\usepackage{amsthm}				% für proof
\usepackage{mathrsfs}			% für \mathscr
\usepackage{latexsym}
\usepackage{marvosym}				% für Lightning

\usepackage{fontspec} 			% funktioniert nur mit den neueren Compilern z.B. XeLaTeX
\usepackage{microtype}			% für bessere Worttrennung
\usepackage[ngerman]{babel} 	% Spracheinstellung
\usepackage{lmodern}			% verändert verwendete Schriftart, damit sie weniger pixelig ist

\usepackage{verbatim}
\usepackage{listings}			% Für Quellcode

\usepackage{graphicx}
\usepackage{tabularx}			% für Tabellen mit gleicher Spaltenbreite und automatischen Umbrüchen
\usepackage{fullpage}
\usepackage{multirow}			% für multirow in tabulars
\usepackage{rotate}
\usepackage[cmyk,table]{xcolor} % um Farben zu benutzen, kann mehr als das Paket color
\usepackage[					% Verlinkungen
	colorlinks,					% farbige Schrift, statt farbiger Rahmen
	linktocpage,				% verlinkt im Abb.Verzeichnis Seitenzahl statt Bildunterschrift
	linkcolor=blue				% setzt Farbe der Links auf blau
	]{hyperref}					% nur für digitale Anwendungen, url = "http://www.example.com"
\usepackage{url}				% für Webadressen wie e-mail usw.: "\url{http://www.example.com}"

\usepackage{enumerate}			% für versch. Aufzählungezeichen wie z.B. a)
\usepackage{xspace}				% folgt ein Leerzeichen nach einem \Befehl, wird es nicht verschluckt.
\usepackage{cancel}				% für das Durchstreichen u.a. in Matheformeln mit \cancel
\usepackage{float}              % zum Forcieren der Position von figure-Umgebungen

% zum Zeichnen (u.a. von Graphen)
\usepackage{fp}
\usepackage{tikz}
\usetikzlibrary{tikzmark}			% für \tikzmark{toRemember}
\usetikzlibrary{positioning}	% verbesserte Positionierung der Knoten
\usetikzlibrary{automata}		% für Automaten (GTI)
\usetikzlibrary{arrows}
\usetikzlibrary{shapes}
\usetikzlibrary{decorations.pathmorphing}
\usetikzlibrary{decorations.pathreplacing}
\usetikzlibrary{decorations.shapes}
\usetikzlibrary{decorations.text}

% //////////////////// Syntaxhighlighting ////////////////////
\lstloadlanguages{Python, Haskell, [LaTeX]TeX, Java}
\lstset{
   basicstyle=\footnotesize\ttfamily,	% \scriptsize the size of the fonts that are used for the code
   backgroundcolor = \color{bgcolour},	% legt Farbe der Box fest
   breakatwhitespace=false,	% sets if automatic breaks should only happen at whitespace
   breaklines=true,			% sets automatic line breaking
   captionpos=t,				% sets the caption-position to bottom, t for top
   commentstyle=\color{codeblue}\ttfamily,% comment style
   frame=single,				% adds a frame around the code
   keepspaces=true,			% keeps spaces in text, useful for keeping indentation
							% of code (possibly needs columns=flexible)
   keywordstyle=\bfseries\ttfamily\color{codepurple},% keyword style
   numbers=left,				% where to put the line-numbers;
   							% possible values are (none, left, right)
   numberstyle=\tiny\color{codegreen},	% the style that is used for the line-numbers
   numbersep=5pt,			% how far the line-numbers are from the code
   stepnumber=1,				% nummeriert nur jede i-te Zeile
   showspaces=false,			% show spaces everywhere adding particular underscores;
							% it overrides 'showstringspaces'
   showstringspaces=false,	% underline spaces within strings only
   showtabs=false,			% show tabs within strings adding particular underscores
   flexiblecolumns=false,
   tabsize=1,				% the step between two line-numbers. If 1: each line will be numbered
   stringstyle=\color{orange}\ttfamily,	% string literal style
   numberblanklines=false,				% leere Zeilen werden nicht mitnummeriert
   xleftmargin=1.2em,					% Abstand zum linken Layoutrand
   xrightmargin=0.4em,					% Abstand zum rechten Layoutrand
   aboveskip=2ex, 
}

\lstdefinestyle{py}{
   language=Python,
}
\lstdefinestyle{hs}{
   language=Haskell,
}
\lstdefinestyle{tex}{
	language=[LaTeX]TeX,
	escapeinside={\%*}{*)},     % if you want to add LaTeX within your code
	texcsstyle=*\bfseries\color{blue},% hervorhebung der tex-Schlüsselwörter
	morekeywords={*,$,\{,\},\[,\],lstinputlisting,includegraphics,
	rowcolor,columncolor,listoffigures,lstlistoflistings,
	subsection,subsubsection,textcolor,tableofcontents,colorbox,
	fcolorbox,definecolor,cellcolor,url,linktocpage,subtitle,
	subject,maketitle,usetikzlibrary,node,path,addbibresource,
	printbibliography},% if you want to add more keywords to the set
     numbers=none,
     numbersep=0pt,
     xleftmargin=0.4em,
}

\lstdefinestyle{java}{
	language=Java,
	extendedchars=true,		% lets you use non-ASCII characters;
   						% for 8-bits encodings only, does not work with UTF-8
}

\lstdefinelanguage[x64]{Assembler}     % add a "x64" dialect of Assembler
   [x86masm]{Assembler} % based on the "x86masm" dialect
   % with these extra keywords:
   {morekeywords={CDQE,CQO,CMPSQ,CMPXCHG16B,JRCXZ,LODSQ,MOVSXD, %
                  POPFQ,PUSHFQ,SCASQ,STOSQ,IRETQ,RDTSCP,SWAPGS, %
                  rax,rdx,rcx,rbx,rsi,rdi,rsp,rbp, %
                  r8,r8d,r8w,r8b,r9,r9d,r9w,r9b}
}					% for 8-bits encodings only, does not work with UTF-8

\lstdefinestyle{c}{
	language=c,
	extendedchars=true,		% for 8-bits encodings only, does not work with UTF-8
}

% //////////////////// eigene Kommandos ////////////////////
\newcommand\FU{Freie Universität Berlin\xspace}% benötigt package xspace
\newcommand\gdw{g.\,d.\,w.\xspace}
\newcommand\oBdA{o.\,B.\,d.\,A.\xspace}
\newcommand{\Eu}{\texteuro}
\newcommand\N{\mathbb{N}\xspace}
\newcommand\Q{\mathbb{Q}\xspace}
\newcommand\R{\mathbb{R}\xspace}
\newcommand\Z{\mathbb{Z}\xspace}
\newcommand\ohneNull{\ensuremath{\backslash\lbrace 0\rbrace}}% \{0}
\let\dhALT\dh	% Schreibt Befehl \dh in \dhALT um
\renewcommand\dh{d.\,h.\xspace}	%renew überschreibt command \dh
\newcommand\Bolt{\;\text{\LARGE\raisebox{-0.3em}{\Lightning}\normalsize}\xspace}% Blitz
\newcommand\zz{\ensuremath{\raisebox{+0.25ex}{Z}% zu zeigen
			\kern-0.4em\raisebox{-0.25ex}{Z}%
			\;\xspace}}
\newcommand{\from}{\ensuremath{\colon}}
\newcommand{\floor}[1]{\lfloor{#1}\rfloor}
\newcommand{\ceil}[1]{\lceil{#1}\rceil}
 \renewcommand{\L}{\ensuremath{\mathcal{L}}\xspace}
 \renewcommand{\P}{\ensuremath{\mathcal{P}}\xspace}
 \newcommand{\NL}{\ensuremath{\mathcal{N}\kern-0.2em\mathcal{L}}\xspace}
 \newcommand{\NP}{\ensuremath{\mathcal{NP}}\xspace}

% //////////////////// Mathefunktionen ////////////////////
\DeclareMathOperator{\Landau}{\mathcal{O}}
\DeclareMathOperator{\True}{True}
\DeclareMathOperator{\False}{False}

% //////////////////// eigene Theoreme ////////////////////
\newtheorem{theorem}{Satz}
\newtheorem{corollary}[theorem]{Folgerung}
\newtheorem{lemma}[theorem]{Lemma}
\newtheorem{observation}[theorem]{Beobachtung}
\newtheorem{definition}[theorem]{Definition}
\newtheorem{Literatur}[theorem]{Literatur}
% konfiguriert proof
\makeatletter
\newenvironment{Proof}[1][\proofname]{\par
  \pushQED{\qed}%
  \normalfont \topsep6\p@\@plus6\p@\relax
  \trivlist
  \item[\hskip\labelsep
%         \itshape
        \bfseries
    #1\@addpunct{.}]\ignorespaces
}{%
  \popQED\endtrivlist\@endpefalse
}
\makeatother

% //////////////////// eigene Farben ////////////////////
\let\definecolor=\xdefinecolor
\definecolor{FUgreen}{RGB}{153,204,0}
\definecolor{FUblue}{RGB}{0,51,102}

\definecolor{middlegray}{rgb}{0.5,0.5,0.5}
\definecolor{lightgray}{rgb}{0.8,0.8,0.8}
\definecolor{orange}{rgb}{0.8,0.3,0.3}
\definecolor{azur}{rgb}{0,0.7,1}
\definecolor{yac}{rgb}{0.6,0.6,0.1}
\definecolor{Pink}{rgb}{1,0,0.6}

\definecolor{bgcolour}{rgb}{0.97,0.97,0.97}
\definecolor{codegreen}{rgb}{0,0.6,0}
\definecolor{codegray}{rgb}{0.35,0.35,0.35}
\definecolor{codepurple}{rgb}{0.58,0,0.82}
\definecolor{codeblue}{rgb}{0.4,0.5,1}

% //////////////////// eigene Settings ////////////////////

\textheight = 230mm		% Höhe des Satzspiegels / Layouts
\footskip = 10ex			% Abstand zw. Fußzeile und Grundlinie letzter Textzeile
\parindent 0pt			% verhindert Einrückung der 1. Zeile eines Absatzes
\setkomafont{sectioning}{\rmfamily\bfseries}% setzt Ü-Schriften in Serifen, {disposition}											% bindet Header ein (WICHTIG)
\usepackage{graphicx}
\usepackage{amsmath}
\usepackage{amssymb}
\usepackage{fancyvrb}

\newcommand{\dozent}{Prof. R. Rojas}					% <-- Names des Dozenten eintragen
\newcommand{\projectNo}{3}
\newcommand{\veranstaltung}{Mustererkennung}
\newcommand{\semester}{WS17/18}
\newcommand{\studenten}{Boyan Hristov, Nedeltscho Petrov}
% /////////////////////// BEGIN DOKUMENT /////////////////////////


\begin{document}
% /////////////////////// BEGIN TITLEPAGE /////////////////////////
\begin{titlepage}
	\subject{\dozent}
	\title{\veranstaltung, \semester}
	\subtitle{\Large Übungsblatt \projectNo\\ \large\vspace{1ex} }
	\author{\studenten}
	\date{\normalsize \today}
\end{titlepage}

\maketitle								% Erstellt das Titelblatt
\vspace*{-9cm}							% rückt Logo an den oberen Seitenrand
\makebox[\dimexpr\textwidth+1cm][r]{	%rechtsbündig und geht rechts 1cm über Layout hinaus
	\includegraphics[width=0.4\textwidth]{src/fu_logo} % fügt FU-Logo ein
}
% /////////////////////// END TITLEPAGE /////////////////////////

\vspace{7cm}							% Abstand
\rule{\linewidth}{0.8pt}				% horizontale Linie										% erstellt die Titelseite


Link zum Git Repository: \url{https://github.com/BoyanH/FU-MachineLearning-17-18/tree/master/Solutions/Homework\projectNo}

% /////////////////////// Aufgabe 1 /////////////////////////

\section*{Klassifikation mit Gauss Verteilung}


\section*{Score (Output des Programs) bzw. Analyse}
Das ist die Ausgabe des Programs und damit auch das Score
\begin{lstlisting}
Score 3 vs 5: 0.923312883436
Score 3 vs 7: 0.929712460064
Score 3 vs 8: 0.843373493976
Score 5 vs 7: 0.931596091205
Score 5 vs 8: 0.874233128834
Score 7 vs 8: 0.837060702875
Score 3 vs 5 vs 7 vs 8: 0.838810641628
Score all: 0.840558046836
Score all with train data: 0.969551501852
\end{lstlisting}

Leider haben wir sehr komisches Score bei der letzten Hausaufgabe bekommen und können das jetzige anhand von selbsterstellte Daten von linearen Regression vergleichen. Mit den Daten aus der Musterlösung, die im Tutorium vorgestellt wurde, wurden wir sagen, dass die jetzige Lösung mit Gauss Verteilung nicht besser (sogar schlechter) ist, als diese mit linearen Regression (da war die Unterscheidung zwischen 5 und 7 mit $94+\%$ Genauigkeit ).

Wenn man alle Ziffer vergleicht, bekommt man ein Score von $84.0558046836\%$ Genauigkeit mit den Testdaten und $96.9551501852\%$ Genauigkeit mit den Traindaten. D.h. wir lernen immer noch ganz viel "auswendig" bzw. es gibt relativ viel "overfitting". 

Das Programm funktioniert ungefähr so schnell wie bei der Klassifikation mit linearen Regression. Die Implementierung war aber schwieriger und kostet mehr Speicher damit es performant ist (Zentren, Konvergenzmatritzen, pseudo-invertierte Konvergenzmatritzen und Determinanten). Bei der Implementierung mit linearen Regression musste man nur die "Hat" Matritzen speichern.

Es ist aber leichter für mehrere Klassen zu implementieren. Man muss nicht separat die Ergebnissen vom mehreren Klassifikatoren vergleichen um zwischen mehrere Klassen zu unterscheiden. Man kann direkt die Wahrscheinlichkeit, dass ein Punkt zu einer Klasse gehört berechnen, und dann direkt die Klasse mit der Höhsten Wahrscheinlichkeit wählen.

\section*{Erklärung interessanteren Teilen des Codes}

Die Fit Methode ist diesmal nicht besonders interessant, da wir dort nur die Zentren berechnen (Alle Punkte aufsummieren und durch ihre Anzahl teilen) und die Konvergenzmatrixberechnungsfunktion aufrufen. Diese wäre aber interessant, da wir die vektorisiert haben und die Hauptimplementierung ausgelagert haben.

Die Summe und die Teilung durch der Anzahl kann man hier auch sehen. Sonst haben wir numpy.vectorize benutzt, um die Berechnung der einzelnen Summaren zu beschleunigen.
\begin{lstlisting}[style=py]
self.covariance_matrix[label] = np.vectorize(GaussianClassifier.covariance_for_point, signature='(m),(n)->(m,m)')(
                points_per_label[label], self.centers[label]).sum(axis=0) / len(points_per_label[label])
\end{lstlisting}

Wir nutzen hier die Formel aus der Vorlesung. Das einzigste Unterschied hier ist, dass wir den ersten Vektor transponieren und nicht der zweiten.
Das ist so, da wir am Ende eine Matrix bekommen wollen und Numpy interpretiert eine 1-dimensionale Liste als eine Matritze der Form 1xn und nicht mx1 (Vektor) als wir
diese betrachten wollen und umgekehrt für das zweite.
\begin{lstlisting}[style=py]
def covariance_for_point(point, center):
        return np.matrix(point - center, dtype=np.float64).T.dot(np.matrix(point - center, dtype=np.float64))
\end{lstlisting}

Bei der Berechnung der Wahrscheinlichkeit haben wir auch die Formel aus der Vorlesung benutzt. Wichtig wäre hier, dass es zu eine Determinante = 0 kommen könnte,
wenn wir linear abhängige Features haben. Das ist oft der Fall bei manchen Ziffern, da wir die Pixels als Features betrachten und diese oft abhängig sind. Z.b wenn eine Ziffer ein Horizontales Strich hat, dann falls ein Pixel grauer ist, sind mit größem Wahrscheinlichkeit auch diese daneben (links und rechts) auch grauer.

Damit wir die 0 Determinante umgehen (wegen Division mit 0), haben wir $2 * \pi * ||X||$ ausgelagert und benutzen später das Minimum von die und eine kleine Delta.
\begin{lstlisting}[style=py]
def get_possibility_for_class(self, point_class, point):
        two_pi_det = 2 * math.pi * self.covariance_matrix_det[point_class]
        left_side = 1 / max(0.2, math.sqrt(two_pi_det))
        right_side = math.e**(-0.5 * (point - self.centers[point_class]).T.
                              dot(self.covariance_matrix_pinv[point_class]).dot(point - self.centers[point_class]))

        return left_side * right_side
\end{lstlisting}

Die Voraussage machen wir, in dem wir die Wahrscheinlichkeit für je Klasse berechnen und dann diese mit höhsten nehmen.
\begin{lstlisting}[style=py]
possibilities = list(map(lambda x: self.get_possibility_for_class(x, point), self.classes))
        winning_index = possibilities.index(max(possibilities))
        return self.classes[winning_index]
\end{lstlisting}

\section*{Vollständiges Code}

\begin{lstlisting}[style]
from Classifier import Classifier
from Parser import *
import numpy as np
import math


class GaussianClassifier(Classifier):
    @staticmethod
    def covariance_for_point(point, center):
        # calculate covariance for a single point (well not really, but a single summar thingy)
        # idea of this method is to easily vectorize it with np.vectorize

        # implementation specific:
        # in the formula, the first vector should be transposed and the second not
        # this is only done, because we need to receive a matrix at the end
        # with the way numpy handles single vectors, we actually need to transpose the first one
        # and not the second in order to do that what we are used to in math
        # (numpy treats a 1-dimensional array as 1xn matrix and not as nx1 as we want to)
        return np.matrix(point - center, dtype=np.float64).T.dot(np.matrix(point - center, dtype=np.float64))

    def __init__(self, train_data, classes = [x for  x in range(10)]):
        """
        :param classes: list of classes the classifier should train itself to distinguish
                        (e.g [3,5] for 3 vs 5 classifier) default is all digits
        :param trainData:
        :param trainLabels:
        :param testData:
        :param testLabels:
        """

        self.centers = {}
        self.covariance_matrix = {}
        self.covariance_matrix_det = {}
        self.covariance_matrix_pinv = {}
        (train_labels, train_points) = get_labels_and_points_from_data(train_data, classes)
        self.classes = classes
        self.fit(train_labels, train_points)

    def fit(self, train_labels, train_points):
        assert(len(train_labels) == len(train_points))
        points_per_label = {}

        # sort points in a dictionary, separated by classes
        # eg {3: [first 256 dimension vector, second 256 dimensional vector, etc.], 5: ... , ...}
        for idx, point in enumerate(train_points):
            current_label = train_labels[idx]
            if current_label not in points_per_label:
                points_per_label[current_label] = [point]
            else:
                points_per_label[current_label].append(point)

        # then for each class, find the centroid and the covariance matrix
        # for optimization reasons, we also save the inverse of the covariance matrix and it's determinant
        for label in points_per_label:
            # average of all points from the current class (with axis 0, so row-wise average)
            self.centers[label] = np.array(points_per_label[label], dtype=np.float64).mean(0)
            # calculate covariance matrix using vectorization (see covariance_for_point static method)
            # using the formula 1/n*(sum_i((point-center)(point-center)T))
            self.covariance_matrix[label] = np.vectorize(GaussianClassifier.covariance_for_point, signature='(m),(n)->(m,m)')(
                points_per_label[label], self.centers[label]).sum(axis=0) / len(points_per_label[label])
            # also calculate and save determinant and pseudo-inverse of matrix for performance reasons
            self.covariance_matrix_det[label] = np.linalg.det(self.covariance_matrix[label])
            self.covariance_matrix_pinv[label] = np.linalg.pinv(self.covariance_matrix[label])

    def predict(self, X):
        return list(map(lambda x: self.predict_single(x), X))

    def predict_single(self, point):
        possibilities = list(map(lambda x: self.get_possibility_for_class(x, point), self.classes))
        winning_index = possibilities.index(max(possibilities))

        return self.classes[winning_index]

    def get_possibility_for_class(self, point_class, point):
        # using the formula from the lecture, calculate the probability for a point with coordinates to
        # be part of a class

        # only important thing here is that 2*pi*det(covariance_matrix) can be zero
        # (in case the covariance_matrix doesn't have a full rank (when we have identical values for some features
        # this can often be the case because of the white pixels at the edges)),
        # so we use
        # np.nextafter to replace any zeros with a reaaaaly small float (because of DivideByZero exceptions...)

        two_pi_det = 2 * math.pi * self.covariance_matrix_det[point_class]
        left_side = 1 / max(0.2, math.sqrt(two_pi_det))
        right_side = math.e**(-0.5 * (point - self.centers[point_class]).T.
                              dot(self.covariance_matrix_pinv[point_class]).dot(point - self.centers[point_class]))

        return left_side * right_side

train_data = parse_data_file('./Dataset/train')
test_data = parse_data_file('./Dataset/test')

three_vs_five = GaussianClassifier(train_data, [3,5])
(three_vs_five_test_labels, three_vs_five_test_data) = get_labels_and_points_from_data(test_data, [3,5])
print("Score 3 vs 5: {}%".format(three_vs_five.score(three_vs_five_test_data, three_vs_five_test_labels)))

three_vs_seven = GaussianClassifier(train_data, [3,7])
(three_vs_seven_test_labels, three_vs_seven_test_data) = get_labels_and_points_from_data(test_data, [3,7])
print("Score 3 vs 7: {}%".format(three_vs_seven.score(three_vs_seven_test_data, three_vs_seven_test_labels)))

three_vs_eight = GaussianClassifier(train_data, [3,8])
(three_vs_eight_test_labels, three_vs_eight_test_data) = get_labels_and_points_from_data(test_data, [3,8])
print("Score 3 vs 8: {}%".format(three_vs_eight.score(three_vs_eight_test_data, three_vs_eight_test_labels)))

five_vs_seven = GaussianClassifier(train_data, [5,7])
(five_vs_seven_test_labels, five_vs_seven_test_data) = get_labels_and_points_from_data(test_data, [5,7])
print("Score 5 vs 7: {}%".format(five_vs_seven.score(five_vs_seven_test_data, five_vs_seven_test_labels)))

five_vs_eight = GaussianClassifier(train_data, [5,8])
(five_vs_eight_test_labels, five_vs_eight_test_data) = get_labels_and_points_from_data(test_data, [5,8])
print("Score 5 vs 8: {}%".format(five_vs_eight.score(five_vs_eight_test_data, five_vs_eight_test_labels)))

seven_vs_eight = GaussianClassifier(train_data, [7,8])
(seven_vs_eight_test_labels, seven_vs_eight_test_data) = get_labels_and_points_from_data(test_data, [7,8])
print("Score 7 vs 8: {}%".format(seven_vs_eight.score(seven_vs_eight_test_data, seven_vs_eight_test_labels)))

combined = GaussianClassifier(train_data, [3, 5, 7, 8])
(combined_test_labels, combined_test_data) = get_labels_and_points_from_data(test_data, [3,5,7,8])
print("Score 3 vs 5 vs 7 vs 8: {}%".format(combined.score(combined_test_data, combined_test_labels)))

all_digits = [x for x in range(10)]
all_classifier = GaussianClassifier(train_data, all_digits)
(all_test_labels, all_test_data) = get_labels_and_points_from_data(test_data, all_digits)
print("Score all: {}%".format(all_classifier.score(all_test_data, all_test_labels)))

all_digits = [x for x in range(10)]
all_classifier = GaussianClassifier(train_data, all_digits)
(all_train_labels, all_train_data) = get_labels_and_points_from_data(train_data, all_digits)
print("Score all: {}".format(all_classifier.score(all_train_labels, all_train_data)))

\end{lstlisting}

\section*{Und das simple Parser}

\begin{lstlisting}[style=py]
import csv
import numpy as np


def parse_data_file(file_name):
    file = open(file_name, 'rt')
    reader = csv.reader(file, delimiter=' ', quoting=csv.QUOTE_NONE)
    data = []

    for row in reader:
        filtered = list(filter(lambda x: x != '', row))
        data.append(list(map(lambda x: float(x), filtered)))

    return data


def get_labels_and_points_from_data(data, classes):
    data = list(filter(lambda x: int(x[0]) in classes, data))
    labels = np.array(list(map(lambda x: int(x[0]), data)))
    points = np.array(list(map(lambda x: x[1:], data)), dtype=np.float64)

    return labels, points

\end{lstlisting}


% /////////////////////// END DOKUMENT /////////////////////////
\end{document}
